\clearpage
\ExposeTOC
\begin{center}
{\Large\bfseries 报告 XXII}\par
\vspace{1.5cm}
{\Large\bfseries zeta 函数的一个同余公式}\par
\vspace{0.8cm}
{\large N. Katz}\par
\end{center}

\vspace{1.5cm}
\section*{0. 引言}

设 $X$ 是 $\mathbb F_q$ 上的固有概形. 将其 zeta 函数视为 $\mathbb Z[[T]]$ 的元素,
即可模 $p$ 约化; 我们希望计算它模 $p$ 的结果.
所得表达式正如从“Woods Hole”型不动点公式所预期的那样:
它是适当的“Frobenius”在下列三种上同调空间中的任一种上作用的特征多项式的交错积:

\begin{enumerate}[label=\arabic*)]
\item 凝聚上同调 $H^i(X,\mathcal O_X)$ (3.1.1);
\item de Rham 上同调
  $H^i_{\mathrm{DR}}(X)=\mathbb H^i(X,\Omega_X^\bullet)$ (3.1.1);
\item 平展上同调 $H^i_{\mathrm{et}}(X,\mathbb Z/p\mathbb Z)$ (3.1.2).
\end{enumerate}

此表达式应视为 Woods Hole 迹公式的一个特例的加强,
即结构层与 Frobenius 自同态的情形. 此特例确实是同余公式的推论 (参见 3.2.1).

本报告的方法完全是朴素的: 先归结到超曲面情形,
再将 Dwork 对超曲面实际 zeta 函数的计算“模 $p$ 约化”.
仔细作此“约化”, 便得到 Ax~[1] 对著名的 Warning--Chevalley 定理~[8] 的加强 5.0.6
(在光滑情形, Dwork 也得到了这一加强 [4, \S 7]).

Deligne 刚刚找到另一种完全上同调的方法, 遵循 SGA~5 的思想,
也适用于“扭系数”, 应会出现在 SGA~5 的新版中.
此法基于 Woods Hole 型迹公式以及 Deligne 的“对称 Künneth 公式”
(SGA~4, XVII~5.5.21).

\section*{1. 回顾“$p$ 线性”运算}

\subsection*{1.0.}
设 $V$ 是特征 $p>0$ 的域 $k$ 上的有限维向量空间, $q$ 是 $p$ 的幂. 映射

\begin{equation}\tag{1.0.1}
\varphi:V\longrightarrow V
\end{equation}

称为 $q$-\emph{线性的}, 如果它是加性的并且满足

\begin{equation}\tag{1.0.2}
\varphi(\lambda v)=\lambda^q\varphi(v)
\qquad\text{对每个 }\lambda\in k,\ v\in V.
\end{equation}

\subsection*{1.0.3.}
注意, $q$ 线性自同态的 $n$ 次迭代 $\varphi^n$ 是 $q^n$ 线性的;
若 $k$ 不完全, 则 $q$ 线性自同态的像 $\varphi(V)$ 未必是 $k$ 向量子空间.

\subsection*{1.0.4.}
对每个扩张 $K/k$, $V$ 的 $q$ 线性自同态 $\varphi$ 唯一地给出
$V_K=V\otimes_kK$ 的一个 $q$ 线性自同态 $\varphi_K$, 满足

\begin{equation}\tag{1.0.5}
\varphi_K(\mu\otimes v)=\mu^q\varphi(v)
\qquad\text{对 }\mu\in K,\ v\in V.
\end{equation}

\subsection*{1.0.6.}
$V$ 的 $q$ 线性自同态 $\varphi$ 称为\emph{半单的 (semisimple)},
如果其像在 $k$ 上张成 $V$; 称为\emph{幂零的}, 如果某次迭代为零.
定义下列 $\varphi$ 稳定的 $k$ 子空间:

\begin{equation}\tag{1.0.7}
V_{\mathrm{ss}}=
\bigcap_{n\geq1}\langle\varphi^n(V)\rangle_k,
\end{equation}

\begin{equation}\tag{1.0.8}
V_{\mathrm{nilp}}=
\bigcup_{n\geq1}\Ker(\varphi^n:V\longrightarrow V).
\end{equation}

显然, $\varphi$ 在 $V_{\mathrm{ss}}$ 上的限制是半单的,
而在 $V_{\mathrm{nilp}}$ 上的限制是幂零的. 若 $k$ 完全, 则

\begin{equation}\tag{1.0.9}
V=V_{\mathrm{ss}}\oplus V_{\mathrm{nilp}}.
\end{equation}

然而若 $k$ 不完全, 则 (1.0.9) 不成立, 因为存在单射的非半单 $q$ 线性自同态.
对任意扩张 $K/k$, 有

\begin{equation}\tag{1.0.10}
(V_K)_{\mathrm{ss}}=V_{\mathrm{ss}}\otimes_kK.
\end{equation}

